\documentclass[11pt]{article}
\usepackage[margin=1in]{geometry}
\usepackage{amsmath,amssymb,amsfonts,mathtools}
\usepackage{array,booktabs,enumitem,graphicx,xcolor,hyperref}
\usepackage[T1]{fontenc}
\usepackage[utf8]{inputenc}
\title{Twisted Black p-Brane Solutions in String Theory}
\author{Ashoke Sen}
\date{}
\begin{document}
\maketitle
\begin{abstract}

It has been shown that given a classical background in string theory which is independent of \$d\$ of the space-time coordinates, we can generate other classical backgrounds by \$O(d)\textbackslash{}otimes O(d)\$ transformation on the solution. We study the effect of this transformation on the known black \$p\$-brane solutions in string theory, and show how these transformations produce new classical solutions labelled by extra continuous parameters and containing background antisymmetric tensor field.

\end{abstract}

\section{Introduction}

It has been shown that given a classical background in string theory which is independent of math expression of the space-time coordinates, we can generate other classical backgrounds by math expression transformation on the solution. We study the effect of this transformation on the known black math expression -brane solutions in string theory, and show how these transformations produce new classical solutions labelled by extra continuous parameters and containing background antisymmetric tensor field.

\section{Method}

math expression where math expression , math expression and math expression denote the graviton, the dilaton, and the antisymmetric tensor fields respectively, math expression + cyclic permutations, math expression denotes the math expression dimensional Ricci scalar, and math expression is the cosmological constant equal to math expression for bosonic string and math expression for fermionic string. (For simplicity we have set the other background massless fields, which appear in fermionic string theories, to zero.) Let us now split the coordinates math expression into two sets math expression and math expression ( math expression , math expression ) and consider backgrounds independent of math expression . Let us further concentrate on backgrounds where math expression , i.e. to backgrounds of the form math expression , math expression . In this case, after an integration by parts, the action \textbackslash{} can be shown to take the form: math expression - d\textasciicircum{}d \& d\textasciicircum{} D-d e\textasciicircum{} -

\begin{equation}
\label{eq:compact}
L(\theta) = \sum_{i=1}^{n} \ell(x_i, \theta)
\end{equation}
Equation~\ref{eq:compact} is included to keep the sample paper-like while remaining small.

\section{Conclusion}

math expression If the coordinates math expression are all of Euclidean signature, this action is invariant under an math expression transformation on math expression , math expression and math expression , given by, math expression M 1 4 M
\end{document}
